\documentclass{article}
\usepackage{amsmath, amssymb, amsthm}
\usepackage{hyperref}
\usepackage{geometry}
\geometry{margin=1in}

\title{Erd\H{o}s Problem \#274}
\date{Last updated: 2025-08-31}
\author{Source: erdosproblems.com}

\begin{document}
\maketitle

\noindent\textbf{Status:} OPEN \quad \textbf{No prize}

\noindent\textbf{Tags:} group theory, covering systems

\medskip
\noindent\textbf{Problem Statement:}

If $G$ is a group then can there exist an exact covering of $G$ by more than one cosets of different sizes? (i.e. each element is contained in exactly one of the cosets)

\bigskip
\noindent\textbf{Remarks:}

A question of Herzog and Sch\"{onheim}, who conjectured more generally that if $G$ is any (not necessarily finite) group and $a_1G_1,\ldots,a_kG_k$ are finitely many cosets of subgroups of $G$ with distinct indices $[G:G_i]$ then the $a_iG_i$ cannot form a partition of $G$.

This conjecture was proved in the case when all the $G_i$ are subnormal in $G$ by Sun \cite{Su04}. In particular if $G$ is abelian (which was the special case asked about in \cite{Er77c} and \cite{ErGr80}) the answer to the original question is no.

Margolis and Schnabel \cite{MaSc19} proved this conjecture for all groups $G$ of size $<1440$.

\bigskip
\noindent\textbf{References:}

[Er77c] Erd\H{o}s, Paul, Problems and results on combinatorial number theory. III. Number theory day (Proc. Conf., Rockefeller Univ.,
New York, 1976) (1977), 43-72.
 
 [ErGr80] Erd\H{o}s, P. and Graham, R., Old and new problems and results in combinatorial number theory. Monographies de L'Enseignement Mathematique (1980).
 
 [MaSc19] Margolis, Leo and Schnabel, Ofir, The {H}erzog-Sch\"onheim conjecture for small groups and
harmonic subgroups. Beitr. Algebra Geom. (2019), 399--418.
 
 [Su04] Sun, Zhi-Wei, On the {H}erzog-Sch\"onheim conjecture for uniform covers of
groups. J. Algebra (2004), 153--175.

\bigskip
\noindent\small{Source: \url{https://www.erdosproblems.com/274}}
\end{document}
